\subsection{Evaluation experimentale controlee des recommandations}

Afin de repondre a la remarque relative aux metriques \textit{Precision@K},
\textit{Recall@K} et \textit{nDCG@K}, une evaluation experimentale controlee
a ete definie. Cette evaluation ne vise pas a demontrer une efficacite
pedagogique generalisable, mais a verifier la coherence des recommandations
produites par le moteur adaptatif par rapport a un oracle pedagogique construit
a partir des regles implementees dans AdaptiveEngine.

Le protocole repose sur un mini-cours compose de cinq concepts :
\textit{Variables}, \textit{Conditions}, \textit{Boucles}, \textit{Fonctions}
et \textit{Tableaux}. Le graphe de prerequis utilise pour cette evaluation est
lineaire : Variables $\rightarrow$ Conditions $\rightarrow$ Boucles
$\rightarrow$ Fonctions $\rightarrow$ Tableaux. Les recommandations attendues
sont definies selon les regles pedagogiques suivantes : recommander le concept
echoue, recommander un prerequis faible lorsqu'il existe, ne pas recommander un
concept avance dont les prerequis ne sont pas maitrises, et privilegier les
concepts proches de la lacune detectee dans le graphe de competences.

Pour chaque scenario, les trois premieres recommandations sont comparees a la
liste des recommandations attendues. Les metriques utilisees sont definies
comme suit :

\[
Precision@K = \frac{\text{nombre de recommandations pertinentes dans le top } K}{K}
\]

\[
Recall@K = \frac{\text{nombre de recommandations pertinentes dans le top } K}
{\text{nombre total de recommandations pertinentes attendues}}
\]

\[
DCG@K = \sum_{i=1}^{K} \frac{rel_i}{\log_2(i+1)}
\]

\[
nDCG@K = \frac{DCG@K}{IDCG@K}
\]

ou $rel_i = 1$ si la recommandation a la position $i$ est pertinente, et
$rel_i = 0$ sinon. Ces metriques sont classiques dans l'evaluation des systemes
de recommandation et du classement d'information
\citep{herlocker2004evaluating,jarvelin2002cumulated}.

\begin{table}[H]
\centering
\caption{Evaluation controlee des recommandations du moteur adaptatif}
\begin{tabular}{p{3.2cm}p{3.5cm}ccc}
\hline
Scenario & Situation pedagogique & Precision@3 & Recall@3 & nDCG@3 \\
\hline
controlled.variables.failure & Echec sur Variables & 0.333 & 1.000 & 1.000 \\
controlled.prerequisite.gap & Echec sur Conditions avec prerequis faible & 0.667 & 1.000 & 1.000 \\
controlled.loop.remediation & Echec sur Boucles & 0.333 & 1.000 & 1.000 \\
controlled.ready.progression & Progression normale vers Boucles & 0.667 & 1.000 & 1.000 \\
controlled.advanced.locked & Remediation sur Fonctions & 0.667 & 1.000 & 1.000 \\
\hline
Moyenne & -- & 0.533 & 1.000 & 1.000 \\
\hline
\end{tabular}
\end{table}

Les resultats montrent que, dans ces scenarios controles, les recommandations
pertinentes apparaissent bien dans les trois premieres positions. La
\textit{Precision@3} moyenne reste limitee par le fait que certaines situations
ne comportent qu'une seule recommandation attendue, alors que le top 3 contient
trois elements. Le \textit{Recall@3} et le \textit{nDCG@3} indiquent que les
elements attendus sont recuperes et places en tete du classement. Ces resultats
doivent etre interpretes prudemment : ils valident la coherence du moteur avec
les regles pedagogiques implementees, mais ne constituent pas une preuve
d'amelioration effective de l'apprentissage chez des apprenants reels.
